% Faithful official CCC2022 article16 full AppendixB.
\section*{B Probabilistic Fine-Grained Reduction for $\mathrm{UTIME}[2^{O(\sqrt{n\log n})}]$}
\paragraph{Theorem 67.} $\mathrm{NQL}\times\mathrm{QLSamp}\subseteq\mathrm{AvgBPQL}\Longrightarrow\mathrm{UTIME}[2^{O(\sqrt{n\log n})}]\subseteq\mathrm{RTIME}[2^{O(\sqrt{n\log n})}]$.
We will need the following fine-grained version of Lemma 27.
\paragraph{Lemma 68.} If $\mathrm{NQL}\times\mathrm{QLSamp}\subseteq\mathrm{AvgBPQL}$, then there exists a polynomial $\tau$ such that the following promise problem is in promiseBPP:
\begin{align*}
\Pi_{\mathrm{YES}}&:=\{(x,1^s,1^t)\mid\mathrm{pK}^t(x)\leq s\},\\
\Pi_{\mathrm{NO}}&:=\{(x,1^s,1^t)\mid\mathrm{pK}^{\widetilde O(t)\cdot\tau(|x|)}(x)>s+\log\tau(t)\}.
\end{align*}
\paragraph{Proof.} Define
\[L:=\{(\mathrm{DP}_k(x;z),w,1^s,1^n)\mid\exists M\in\{0,1\}^s,M(w)\text{ outputs }x\in\{0,1\}^n\text{ within }|w|\text{ steps}\},\]
where $c>0$ is some large constant so that $L\in\mathrm{NQL}$. Define a distribution family $\mathcal D:=\{D_{\langle nk+k,t,s,n\rangle}\}$, each member of which does the following: sample $u\sim U_{nk+k}$ and $w\sim U_t$ and then output $(u,w,1^s,1^n)$. (Again, for the simplicity of presentation, we omit the padding, which is of length at most $\widetilde O(t)\cdot\operatorname{poly}(nk)$.) By assumption, $(L,\mathcal D)\in\mathrm{AvgBPQL}$. Let $B$ be a randomized heuristic algorithm for $(L,\mathcal D)$ as described in Lemma 6. Now, define an algorithm $B'$:
\begin{itemize}
\item On input $(x,1^s,1^t)$ with $x\in\{0,1\}^n$, set $k=s+10$, sample $z\sim U_{nk}$ and $w\sim U_t$, and then output $B(\mathrm{DP}_k(x;z),w,1^s,1^n)$.
\end{itemize}
Note that $B'$ runs in time $t_D:=\widetilde O(t)\cdot p(n)$ for some polynomial $p$. Below, we show that $B'$ solves the mentioned promise problem correctly with high probability in the worst case.
First, consider the case that $(x,1^s,1^t)\in\Pi_{\mathrm{YES}}$. By the definitions of $L$ and pK, for any choice of $z\in\{0,1\}^{nk}$,
\[\Pr_w[(\mathrm{DP}_k(x;z),w,1^s,1^n)\in L]\geq2/3.\]
The definition of $B$ then implies that
\[\Pr_{w,z,r_B}[B(\mathrm{DP}_k(x;z),w,1^s,1^n)=1]>1/2,\]
and so
\[\Pr_{r_{B'}}(B'(x,1^s,1^t)=1)>1/2,\tag{42}\]
where $r_B$ denotes the internal randomness of $B$, and $r_{B'}=(w,z,r_B)$ that of $B'$.
Now consider the case that $(x,1^s,1^t)\in\Pi_{\mathrm{NO}}$. For a contradiction, suppose that
\[\Pr_{w,z,r_B}[B(\mathrm{DP}_k(x;z),w,1^s,1^n)=1]>1/3.\tag{43}\]
Recall $k=s+10$. By a counting argument, for randomly selected $u$ and $w$,
\[\Pr_{u,w}[(u,w,1^s,1^n)\in L]\leq\frac{2^s\cdot2^{nk}\cdot2^t}{2^{nk+k+t}}=1/10.\]
Then by definition of $B$,
\[\Pr_{u,w,r_B}[B(u,w,1^s,1^n)=1]\leq1/9.\tag{44}\]
Comparing Equations (43) and (44), we see that $B(-,U_t,1^s,1^n)$ $(2/9)$-distinguishes the distribution $\mathrm{DP}_k(x;U_{nk})$ from uniform. Lemma 22 implies that
\[\mathrm{pK}^{\widetilde O(t)\cdot\tau_0(n)}(x)\leq k+\log\tau_0(t)=s+10+\log\tau_0(t),\]
for some polynomial $\tau_0$ that depends on the polynomials $p$ and $p_{\mathrm{DP}}$ from Lemma 22. This means that $(x,1^s,1^t)$ is not in $\Pi_{\mathrm{NO}}$, by setting $\tau$ properly. This gives the desired contradiction. By definition of $B'$, we have that
\[\Pr_{r_{B'}}(B'(x,1^s,1^t)=1)\leq1/3.\tag{45}\]
By Equations (13) and (16), $B'$ yields a promiseBPP algorithm for $(\Pi_{\mathrm{YES}},\Pi_{\mathrm{NO}})$ via standard error reduction techniques.\hfill$\triangleleft$
The following is a simple corollary of Lemma 47 and Lemma 20.
\paragraph{Lemma 69 (Fine-Grained Weak Symmetry of Information for $\mathrm{pK}^t$).} If $\mathrm{NQL}\times\mathrm{QLSamp}\subseteq\mathrm{AvgBPQL}$, then there exist polynomials $p_0$ and $p$ such that for every sufficiently large $x\in\{0,1\}^n$, $m\in\mathbb N$ and $p_0(n,m)\leq t\leq2^{n+m}$,
\[\Pr_{u\sim\{0,1\}^m}[\mathrm{pK}^t(x\circ u)>\mathrm{pK}^{t\cdot p(nm)}(x)+m-\log p(t)]\geq0.99.\]
We will use the following black-box hitting set generator whose outputs can be small relative to the ``hard string'' used in its construction.
\paragraph{Lemma 70 ([10, Theorem 4.3]).} For any $T,m\in\mathbb N$ with $m\leq2T$, there exists a function $H_{T,m}:\{0,1\}^T\times\{0,1\}^d\to\{0,1\}^m$ and a deterministic procedure $\mathrm{Recon}^{(-)}:\{0,1\}^a\to\{0,1\}^T$, where $d=O(\log T+\log^3m)$ and $a=2m+O(\log T+\log^3m)$ such that, for any $x\in\{0,1\}^T$ and any function $D:\{0,1\}^m\to\{0,1\}$ that $0.1$-avoids $H_{T,m}(x,-)$, there exists advice $\alpha\in\{0,1\}^a$ such that $\mathrm{Recon}^D(\alpha)=x$. Moreover, $H_{T,m}$ can be computed in time $\operatorname{poly}(T)$ and $\mathrm{Recon}^D$ can be computed in time $\operatorname{poly}(T)$ with oracle access to $D$.
We are now ready to show Theorem 67.
\paragraph{Proof of Theorem 67.} Let $L\in\mathrm{UTIME}[T(n)]$ with verifier $V$, where $T(n)=2^{O(\sqrt{n\log n})}$. Fix an input $x\in\{0,1\}^n$. Let $y_x\in\{0,1\}^{T(n)}$ be the unique $L$-witness for $x$. Let $B$ be a promiseBPP algorithm that solves the promise problem $(\Pi_{\mathrm{YES}},\Pi_{\mathrm{NO}})$ from Lemma 68.
Let $t$ be any number such that $q_0(T(n))\leq t\leq2^{n/2}$, where $q_0$ is some polynomial specified later. Let $H:=H_{T(n),m}$ be the hitting set generator from Lemma 70, where $m:=m_x(t)\leq O(n)$ is defined later, and let $d:=O(\log T(n)+\log^3m)$ be the seed length of $H$.
Using $H$ and $B$, we will argue that the time-bounded pK complexity of $y_x$ is small given $x$. First note the following.
\paragraph{Claim 71.} There exists a polynomial $q$ such that for every $q_0(T(n))\leq t\leq2^{n/2}$ and every $z_H\in\{0,1\}^d$,
\[\mathrm{pK}^{t\cdot q(nm)}(x\circ H(y_x;z_H))\leq\mathrm{pK}^t(x)+d+\log q(t).\]
\paragraph{Proof of Claim 71.} Let $L$ be the following language
\begin{multline*}
L':=\{(\mathrm{DP}_k(x\circ H(y;z');z),w01^{T(n)^c},1^s,1^n,1^m)\mid\exists M\in\{0,1\}^s,\\
M(w)\text{ outputs }x\in\{0,1\}^n\text{ within }|w|\text{ steps, and }V(x,y)=1\},
\end{multline*}
where $c>0$ is some constant so that $L'\in\mathrm{NQL}$. Define a distribution family $\mathcal D:=\{D_{\langle(n+m)k+k,t+T(n)^c+1,s,n,m\rangle}\}$, each member of which samples $u\sim U_{(n+m)k+k}$, $w\sim U_t$ and outputs $(u,w01^{T(n)^c},1^s,1^n,1^m)$.
By assumption, $(L',\mathcal D)\in\mathrm{AvgBPQL}$. Let $B'$ be a randomized heuristic algorithm for $(L',\mathcal D)$ as described in Lemma 6.
Define $s:=\mathrm{pK}^t(x)$. By definition of $\mathrm{pK}^t$ and $L'$, for any fixed $z\in\{0,1\}^{(n+m)k}$,
\[\Pr_w[(\mathrm{DP}_k(x\circ H(y_x;z_H);z),w01^{T(n)^c},1^s,1^n,1^m)\in L']\geq2/3.\]
Then by definition of $B'$, it is easy to see that
\[\Pr_{z,w,r_{B'}}[B'(\mathrm{DP}_k(x\circ H(y_x;z_H);z),w01^{T(n)^c},1^s,1^n,1^m)=1]\geq1/2.\tag{46}\]
On the other hand, by a counting argument it is easy to see that for $u$ and $w$ selected uniformly at random, we have
\[\Pr_{u,w}[(u,w01^{T(n)^c},1^s,1^n,1^m)\in L']\leq\frac{2^s\cdot2^t\cdot2^d\cdot2^{(n+m)k}}{2^{|u|}\cdot2^{|w|}}\leq1/n,\]
where the last inequality holds by setting $k:=s+d+\log n$. By the definition of $B'$, we have
\[\Pr_{u,w,r_{B'}}[B'(u,w01^{T(n)^c},1^s,1^n,1^m)=1]\leq o(1).\tag{47}\]
Given Equations (46) and (47), it is clear that $D(-):=B'(-,U_t01^{T(n)^c},1^s,1^n,1^m)$ is a randomized distinguisher for $\mathrm{DP}_k(x\circ H(y_x;z_H);U_{(n+m)k})$. Note that provided that $q_0$ is a large enough polynomial, $D$ needs only $O(\log t)$ bits as advice, and then it runs in time
\[t_D:=\widetilde O(t+T(n)^c)\cdot\operatorname{poly}(nm)\leq\widetilde O(t)\cdot\operatorname{poly}(nm).\]
Lemma 22 implies that
\[\mathrm{pK}^{\widetilde O(t)\cdot q(nm)}(x\circ H(y_x;z_H))\leq k+O(\log t)\leq\mathrm{pK}^t(x)+d+\log q(t),\]
for some large polynomial $q$. This completes the proof of Claim 71.\hfill$\triangleleft$
By Lemma 69, there exist polynomial $p_0$ and $p$ such that for every $p_0(n,m)\leq t_1\leq2^{n+m}$
\[\Pr_{u\sim\{0,1\}^m}[\mathrm{pK}^{t_1}(x\circ u)>\mathrm{pK}^{t_1\cdot p(nm)}(x)+m-\log p(t_1)]\geq0.99.\tag{48}\]
We set $t_1:=\widetilde O(t\cdot q(nm))\cdot\tau(n+m)$, where $\tau$ is the function from Lemma 68. Note that we can make the polynomial $q_0$ to be large enough so that $t_1$ satisfies the condition for which Equation (48) holds. Now we choose $m$ so that
\[\mathrm{pK}^{t_1\cdot p(nm)}(x)+m-\log p(t_1)\geq s+\log\tau(t\cdot q(nm)).\tag{49}\]
That is, we set
\[m:=(\mathrm{pK}^t(x)-\mathrm{pK}^{t\cdot n^c}(x))+c\cdot\log t,\]
where $c>0$ is some large constant, and we use the fact that $t\leq2^{n/2}$ to simplify $t_1$. Define the following (probabilistic) function:
\[D(-):=\neg B(x\circ-,1^s,1^{t\cdot q(nm)}).\]
Observe the following properties of $D$.
\begin{itemize}
\item For every $z\in\{0,1\}^d$, $\Pr_D[D(H(y_x;z))=0]\geq2/3$. This follows from Claim 71 and the correctness of $B$ on solving the promise problem $(\Pi_{\mathrm{YES}},\Pi_{\mathrm{NO}})$ from Lemma 68.
\item For at least $0.99$ fraction of $\beta\in\{0,1\}^m$, $\Pr_D[D(\beta)=1]\geq2/3$. This follows from Equations (48) and (49) (and the correctness of $B$).
\item $D$ runs in time $\operatorname{poly}(t)$.
\end{itemize}
At this point, we can see that $D$ is a randomized function that distinguishes $H(y_x;U_d)$ from $U_m$, and by randomly fixing its internal randomness we get a deterministic distinguisher with good probability. However, in order to utilize the reconstruction procedure in Lemma 70 to recover $y_x$, we need a deterministic function that avoids $H$. By amplifying the success probability of $D$ on the ``promised'' inputs, we can show that randomly fixing the internal randomness of $D$ gives a function that avoids $H$ with high probability. More specifically, using standard error reduction techniques, we can make both the probability in the first two items above at least $1-1/(4\cdot2^m)$, which only causes a multiplicative overhead of $\operatorname{poly}(m)$ in the running time. By abusing notation, we use the same $D$ to denote this new function with small error. Now consider the set $S$ of inputs for which $D$ has this small error. That is,
\[S:=\{H(y_x;z)\mid z\in\{0,1\}^d\}\cup\{\beta\in\{0,1\}^m\mid\Pr_D[D(\beta)=1]\geq1-1/(4\cdot2^m)\}.\]
Note that $|S|\leq2^m$. For $\beta\in S$, let us say that the correct answer of $\beta$ is $0$ (resp. $1$) if $\beta$ is in the first (resp. second) subset in the definition of $S$. By a union bound, we have
\[\Pr_{r_D\sim\{0,1\}^{\operatorname{poly}(t)}}[\exists\beta\in S\text{ such that }D(\beta;r_D)\text{ is not correct}]\leq|S|\cdot\frac1{4\cdot2^m}\leq1/4,\tag{50}\]
where $r_D$ above denotes the internal randomness of $D$. This means with probability at least $2/3$ over a uniformly random $r_D$, $D(-;r_D)$ is a deterministic function that is correct on all the inputs in $S$. In particular, such a function $0.1$-avoids $H$. Let us say that $r_D$ is good if this is true. Let $\mathrm{Recon}^{(-)}:\{0,1\}^a\to\{0,1\}^{T(n)}$ be the reconstruction procedure in Lemma 70, where $a:=2m+O(\log T+\log^3m)$. We have
\begin{align*}
&\Pr_{r_D\sim\{0,1\}^{\operatorname{poly}(t)}}[\exists\alpha\in\{0,1\}^a\text{ such that }\mathrm{Recon}^{D(-;r_D)}(\alpha)=x]\\
&\geq\Pr_{r_D}[\exists\alpha\in\{0,1\}^a\text{ such that }\mathrm{Recon}^{D(-,r_D)}(\alpha)=x\mid r_D\text{ is good}]\cdot\Pr_{r_D}[r_D\text{ is good}]\\
&\geq2/3.\qquad\text{(Lemma70 and Equation50)}
\end{align*}
Note that given $x$, $D$ can be constructed using the integers $s,t,n,m$ and the code for $B$, which can be encoded using $O(\log t)$ bits. Then the above implies
\[\mathrm{pK}^{\operatorname{poly}(t\cdot T)}(y_x\mid x)\leq a+O(\log t)\leq2m+O(\log t+\log^3m).\tag{51}\]
Note that the above holds for every $q_0(T(n))\leq t\leq2^{n/2}$. Also, recall that $m=\mathrm{pK}^t(x)-\mathrm{pK}^{t\cdot n^c}(x)+c\cdot\log t$ for some constant $c>0$. By Lemma 46, there exists some $t^*$ such that
\[q_0(T(n))\leq t^*\leq q_0(T(n))\cdot2^{c\cdot\sqrt{n\log n}}\leq2^{O(\sqrt{n\log n})},\]
and
\[\mathrm{pK}^{t^*}(x)-\mathrm{pK}^{t^*\cdot n^c}(x)\leq O(\sqrt{n\log n}).\]
Then for such $t^*$, we have $m\leq O(\sqrt{n\log n})$. Therefore, plugging this $t^*$ into Equation (51) for $t$, we have that there exists some constant $d>0$ such that
\[\mathrm{pK}^{2^{d\cdot\sqrt{n\log n}}}(y_x\mid x)\leq d\cdot\sqrt{n\log n}.\]
This suggests the following probabilistic algorithm $A$ for solving $L$:
\begin{itemize}
\item On input $x\in\{0,1\}^n$, $A$ samples $w\sim\{0,1\}^{d\cdot\sqrt{n\log n}}$. It then exhaustively searches over all $M\in\{0,1\}^{d\cdot\sqrt{n\log n}}$, running $M(w,x)$ for $2^{d\cdot\sqrt{n\log n}}$ steps to produce some output $y$. $A$ accepts iff a $y$ is obtained such that $V(x,y)=1$.
\end{itemize}
This completes the proof.\hfill$\triangleleft$
